\documentclass[10pt,a4paper]{amsart}
\usepackage{amsmath,amssymb,amsthm,graphicx}
% Retained minimal preamble for the Stacks Project chapter more-algebra.tex (More on Algebra).
% The chapter file pinned-original-source.tex carries only \input{preamble} on its first line, so no
% preamble text lives inside the chapter itself.  The Stacks Project uses one single shared preamble for
% every chapter; the snapshot stages it identically as sources/stacks-chow/preamble.tex and
% sources/stacks-etale-cohomology/preamble.tex.  The block below is copied mechanically, line for line,
% from sources/stacks-chow/preamble.tex lines 14-18 (the Xy-pic setup) and 154-217 (the native theorem
% declarations together with the mathematical macros).  Nothing was added, removed or redefined; the
% wrapper class and the amsmath/amssymb/amsthm/graphicx base packages are supplied by the assembler.
% Chapter: more-algebra.tex -- More on Algebra.
% For commutative diagrams we use Xy-pic
\usepackage[all]{xy}

% We use 2cell for 2-commutative diagrams.
\xyoption{2cell}

% Theorem environments.
%
\theoremstyle{plain}
\newtheorem{theorem}[subsection]{Theorem}
\newtheorem{proposition}[subsection]{Proposition}
\newtheorem{lemma}[subsection]{Lemma}

\theoremstyle{definition}
\newtheorem{definition}[subsection]{Definition}
\newtheorem{example}[subsection]{Example}
\newtheorem{exercise}[subsection]{Exercise}
\newtheorem{situation}[subsection]{Situation}

\theoremstyle{remark}
\newtheorem{remark}[subsection]{Remark}
\newtheorem{remarks}[subsection]{Remarks}

\numberwithin{equation}{subsection}

% Macros
%
\def\lim{\mathop{\mathrm{lim}}\nolimits}
\def\colim{\mathop{\mathrm{colim}}\nolimits}
\def\Spec{\mathop{\mathrm{Spec}}}
\def\Hom{\mathop{\mathrm{Hom}}\nolimits}
\def\Ext{\mathop{\mathrm{Ext}}\nolimits}
\def\SheafHom{\mathop{\mathcal{H}\!\mathit{om}}\nolimits}
\def\SheafExt{\mathop{\mathcal{E}\!\mathit{xt}}\nolimits}
\def\Sch{\mathit{Sch}}
\def\Mor{\mathop{\mathrm{Mor}}\nolimits}
\def\Ob{\mathop{\mathrm{Ob}}\nolimits}
\def\Sh{\mathop{\mathit{Sh}}\nolimits}
\def\NL{\mathop{N\!L}\nolimits}
\def\CH{\mathop{\mathrm{CH}}\nolimits}
\def\proetale{{pro\text{-}\acute{e}tale}}
\def\etale{{\acute{e}tale}}
\def\QCoh{\mathit{QCoh}}
\def\Ker{\mathop{\mathrm{Ker}}}
\def\Im{\mathop{\mathrm{Im}}}
\def\Coker{\mathop{\mathrm{Coker}}}
\def\Coim{\mathop{\mathrm{Coim}}}

% Boxtimes
%
\DeclareMathSymbol{\boxtimes}{\mathbin}{AMSa}{"02}

%
% Macros for moduli stacks/spaces
%
\def\QCohstack{\mathcal{QC}\!\mathit{oh}}
\def\Cohstack{\mathcal{C}\!\mathit{oh}}
\def\Spacesstack{\mathcal{S}\!\mathit{paces}}
\def\Quotfunctor{\mathrm{Quot}}
\def\Hilbfunctor{\mathrm{Hilb}}
\def\Curvesstack{\mathcal{C}\!\mathit{urves}}
\def\Polarizedstack{\mathcal{P}\!\mathit{olarized}}
\def\Complexesstack{\mathcal{C}\!\mathit{omplexes}}
% \Pic is the operator that assigns to X its picard group, usage \Pic(X)
% \Picardstack_{X/B} denotes the Picard stack of X over B
% \Picardfunctor_{X/B} denotes the Picard functor of X over B
\def\Pic{\mathop{\mathrm{Pic}}\nolimits}
\def\Picardstack{\mathcal{P}\!\mathit{ic}}
\def\Picardfunctor{\mathrm{Pic}}
\def\Deformationcategory{\mathcal{D}\!\mathit{ef}}

% --- appended 2026-10-12 (Chief): the project own definitions of the reference/slogan/history
% display environments, copied mechanically from the shared Stacks preamble lines 8-12
% (\usepackage{verbatim} plus the three \newenvironment lines, which the chapter preamble keeps
% verbatim from the project preamble). Nothing added, removed or redefined: these environments
% carry no mathematical content and the project itself typesets them as comments. Without these
% three lines a counted statement whose own text uses them fails with Environment slogan undefined.
\usepackage{verbatim}
\newenvironment{reference}{\comment}{\endcomment}
\newenvironment{slogan}{\comment}{\endcomment}
\newenvironment{history}{\comment}{\endcomment}
\usepackage{cleveref}
\crefname{originallemma}{original lemma}{original lemmas}
\Crefname{originallemma}{Original lemma}{Original lemmas}
\makeatletter
\newlabel{more-algebra-lemma-lift-fs}{{07NR}{source}{}{sourcepointer-1}{}}
\newlabel{more-algebra-lemma-lift-fs@cref}{{[originallemma][0][]07NR}{[1][0][]source}{}{}{}}
\makeatother
\crefname{originallemma}{original lemma}{original lemmas}
\Crefname{originallemma}{Original lemma}{Original lemmas}
\makeatletter
\newlabel{lemma-embed-map-Noetherian-complete-local-rings}{{07NN}{source}{}{sourcepointer-2}{}}
\newlabel{lemma-embed-map-Noetherian-complete-local-rings@cref}{{[originallemma][0][]07NN}{[1][0][]source}{}{}{}}
\makeatother
\crefname{originallemma}{original lemma}{original lemmas}
\Crefname{originallemma}{Original lemma}{Original lemmas}
\makeatletter
\newlabel{algebra-lemma-regular-quotient-regular}{{00NR}{source}{}{sourcepointer-3}{}}
\newlabel{algebra-lemma-regular-quotient-regular@cref}{{[originallemma][0][]00NR}{[1][0][]source}{}{}{}}
\makeatother
\crefname{originallemma}{original lemma}{original lemmas}
\Crefname{originallemma}{Original lemma}{Original lemmas}
\makeatletter
\newlabel{lemma-quotient-power-series-ring-over-Cohen}{{07NM}{source}{}{sourcepointer-4}{}}
\newlabel{lemma-quotient-power-series-ring-over-Cohen@cref}{{[originallemma][0][]07NM}{[1][0][]source}{}{}{}}
\makeatother
\crefname{originallemma}{original lemma}{original lemmas}
\Crefname{originallemma}{Original lemma}{Original lemmas}
\makeatletter
\newlabel{proposition-fs-flat-fibre-fs}{{07NQ}{source}{}{sourcepointer-5}{}}
\newlabel{proposition-fs-flat-fibre-fs@cref}{{[originallemma][0][]07NQ}{[1][0][]source}{}{}{}}
\makeatother
\crefname{originallemma}{original lemma}{original lemmas}
\Crefname{originallemma}{Original lemma}{Original lemmas}
\makeatletter
\newlabel{lemma-base-change-fs}{{07EG}{source}{}{sourcepointer-6}{}}
\newlabel{lemma-base-change-fs@cref}{{[originallemma][0][]07EG}{[1][0][]source}{}{}{}}
\makeatother
\title{Lifting a formally smooth complete local algebra along a residue field extension}
\author{The Stacks Project Authors}
\date{stacks-project:more-algebra.tex:local-snapshot-20261012; GNU FDL 1.2 or later}
\begin{document}
\maketitle
\setcounter{subsection}{1}
\noindent\textit{Editorial note.} All delivered mathematical text is the author's own, copied from the pinned Stacks Project chapter file more-algebra.tex; no statement or proof step was rewritten, reordered or removed.\par\medskip
\noindent\textit{Editorial note.} The chapter file pinned-original-source.tex carries only \textbackslash{}input\{preamble\} on its first line, so no preamble text lives inside the chapter. The retained file sources/stacks-more-algebra/retained-preamble.tex is a mechanical, line-for-line copy of the project's single shared preamble as staged in the snapshot (Xy-pic setup plus theorem declarations and mathematical macros). Nothing was added, removed or redefined, and the excerpt serves only as preamble, never as a span coordinate.\par\medskip
\noindent\textit{Editorial note.} Both counted spans start at the \textbackslash{}begin line of their environment, so the statement and the author's proof are delivered whole; surrounding prose of the section is deliberately outside the counted spans and is carried as standard prerequisites instead.\par\medskip
\noindent\textit{Editorial note.} The citation fields (source\_id, source\_version, authors, source\_url, license.name, source\_title) are identical to those of every other card this lane has produced from this chapter, so the batch selector treats them as one source.\par\medskip
\noindent\textit{Allowed standard prerequisites (editorial).}\begin{itemize}
\item Existence of the diagram of Noetherian complete local rings with a power series ring over a Cohen ring at the top. Lemma \textbackslash{}ref\{lemma-embed-map-Noetherian-complete-local-rings\}, cited at proof lines 10266 and 10295. Provides the ambient power series ring S over R whose special fibre is the given algebra.
\item The kernel of a surjection of regular local rings is generated by part of a regular system of parameters. Algebra, Lemma \textbackslash{}ref\{algebra-lemma-regular-quotient-regular\}, cited at proof lines 10277-10279. Identifies generators x-bar\_1, ..., x-bar\_r of the kernel that extend to a regular system of parameters.
\item The quotient of a power series ring by part of a regular system of parameters is a power series ring over a field or a Cohen ring. Lemma \textbackslash{}ref\{lemma-quotient-power-series-ring-over-Cohen\}, cited at proof lines 10288-10290. Identifies the cut ring as a power series ring, so that the flat-fibre criterion applies to it.
\item Formal smoothness of a local map whose special fibre is formally smooth, and base change of formal smoothness. Proposition \textbackslash{}ref\{proposition-fs-flat-fibre-fs\} and Lemma \textbackslash{}ref\{lemma-base-change-fs\}, cited at proof lines 10300-10304. Give formal smoothness of R to S once the special fibres agree, and of the base change C = S tensor A over A.
\end{itemize}\medskip
\ifdefined\hypertarget\hypertarget{sourcepointer-1}{}\fi
\noindent\textit{Source reference.} Original lemma 07NR: focused-seeds/AGP/sources/stacks-snapshot-20261011/tags.txt, source lines 9771--9771. The lemma's own canonical tag, quoted so that readers can resolve the label inside the delivered text; it points at the counted statement itself and carries no proof content.\par\medskip
\ifdefined\hypertarget\hypertarget{sourcepointer-2}{}\fi
\noindent\textit{Source reference.} Original lemma 07NN: focused-seeds/AGP/sources/stacks-snapshot-20261011/tags.txt, source lines 9768--9768. Navigation pointer to a statement the counted proof refers to but does not contain; nothing load-bearing is delegated to it.\par\medskip
\ifdefined\hypertarget\hypertarget{sourcepointer-3}{}\fi
\noindent\textit{Source reference.} Original lemma 00NR: focused-seeds/AGP/sources/stacks-snapshot-20261011/tags.txt, source lines 869--869. Navigation pointer to a statement the counted proof refers to but does not contain; nothing load-bearing is delegated to it.\par\medskip
\ifdefined\hypertarget\hypertarget{sourcepointer-4}{}\fi
\noindent\textit{Source reference.} Original lemma 07NM: focused-seeds/AGP/sources/stacks-snapshot-20261011/tags.txt, source lines 9767--9767. Navigation pointer to a statement the counted proof refers to but does not contain; nothing load-bearing is delegated to it.\par\medskip
\ifdefined\hypertarget\hypertarget{sourcepointer-5}{}\fi
\noindent\textit{Source reference.} Original lemma 07NQ: focused-seeds/AGP/sources/stacks-snapshot-20261011/tags.txt, source lines 9770--9770. Navigation pointer to a statement the counted proof refers to but does not contain; nothing load-bearing is delegated to it.\par\medskip
\ifdefined\hypertarget\hypertarget{sourcepointer-6}{}\fi
\noindent\textit{Source reference.} Original lemma 07EG: focused-seeds/AGP/sources/stacks-snapshot-20261011/tags.txt, source lines 9446--9446. Navigation pointer to a statement the counted proof refers to but does not contain; nothing load-bearing is delegated to it.\par\medskip
\begin{lemma}
\label{lemma-lift-fs}
Let $A$ be a Noetherian complete local ring with residue field $k$.
Let $B$ be a Noetherian complete local $k$-algebra. Assume $k \to B$
is formally smooth in the $\mathfrak m_B$-adic topology.
Then there exists a Noetherian complete local ring $C$
and a local homomorphism $A \to C$ which is formally smooth
in the $\mathfrak m_C$-adic topology such that $C \otimes_A k \cong B$.
\end{lemma}

\begin{proof}
Choose a diagram
$$
\xymatrix{
S \ar[r] & B \\
R \ar[u] \ar[r] & A \ar[u]
}
$$
as in Lemma \ref{lemma-embed-map-Noetherian-complete-local-rings}.
Let $t_1, \ldots, t_d$ be a regular system of parameters for $R$
with $t_1 = p$ in case the characteristic of $k$ is $p > 0$.
As $B$ and $\overline{S} = S \otimes_R k$
are regular we see that $\Ker(\overline{S} \to B)$ is generated by
elements $\overline{x}_1, \ldots, \overline{x}_r$ which form part of a
regular system of parameters of $\overline{S}$, see
Algebra, Lemma \ref{algebra-lemma-regular-quotient-regular}.
Lift these elements to $x_1, \ldots, x_r \in S$. Then
$t_1, \ldots, t_d, x_1, \ldots, x_r$ is part of a regular system of
parameters for $S$. Hence $S/(x_1, \ldots, x_r)$ is a power
series ring over a field (if the characteristic of $k$ is zero)
or a power series ring over a Cohen ring (if the characteristic of
$k$ is $p > 0$), see
Lemma \ref{lemma-quotient-power-series-ring-over-Cohen}.
Moreover, it is still the case that $R \to S/(x_1, \ldots, x_r)$
maps $t_1, \ldots, t_d$ to a part of a regular system of parameters
of $S/(x_1, \ldots, x_r)$. In other words, we may replace $S$ by
$S/(x_1, \ldots, x_r)$ and assume we have a diagram
$$
\xymatrix{
S \ar[r] & B \\
R \ar[u] \ar[r] & A \ar[u]
}
$$
as in Lemma \ref{lemma-embed-map-Noetherian-complete-local-rings}
with moreover $\overline{S} = B$. In this case $R \to S$ is
formally smooth in the $\mathfrak m_S$-adic topology by
Proposition \ref{proposition-fs-flat-fibre-fs}.
Hence the base change $C = S \otimes_R A$ is formally smooth
over $A$ in the $\mathfrak m_C$-adic topology by
Lemma \ref{lemma-base-change-fs}.
\end{proof}

\end{document}
